\section{Residual terms and integration along paths}
\label{res:section}

We work with a finite tile collection and a fixed angular translation.
Fix a shift $\sigma$ and a parameter $\zeta$ whose real part is the
corresponding vertex: $\Re\zeta_\sigma=\mu$, all other real parts vanish,
and $\norm{\Im\zeta}_2\le L^2$.  All estimates in this section are
uniform in these choices, the number of tiles, the angular depth, and
$\dim\mathcal H$.  We take $L$ sufficiently large that
\begin{equation}
  p>\alpha,\qquad 0<\mu\le\frac1{12}.
  \label{res:small-mu}
\end{equation}
We retain the stopping threshold $M=L^{A_M}$ from
\eqref{hier:mass-cutoff}; its fixed exponent will be chosen at the end
of the residual estimates.

Fix a path top $R$.  We retain the block notation of the hierarchy:
$J_j=v_{i_j}(R)$, $n_j=n_{J_j}$, and $A_j=A_{J_j}$.
The last possible block is the last one whose starting node contains
$R$ as an eligible top.  At this last block there is no path child,
even if other tops continue; the next path mass is defined to be zero.
At every other block write $J_{j+1}$ for the path child and set
\[
  h_j^*=n_j-n_{j+1},\qquad u_j=\frac{h_j^*}{n_j}.
\]
An index is active if $n_j>M$.  The active indices form a prefix.
On $Z$, the mass bound \eqref{eq:src-12} and monotonicity give
\begin{equation}
  \sum_{j\text{ active}}u_j\le C L.
  \label{res:relative-drop-sum}
\end{equation}
Indeed, except for the last active index, use
$1-n_{j+1}/n_j\le\log(n_j/n_{j+1})$ and telescope down to a mass
larger than $M$.  The final summand is at most one.  This proof also
applies when the path stops by ineligibility.

\subsection{Marked residuals}

For $n>0$ and $A\in\mathcal H$, put
\[
  V(n,A)=n+n^{\alpha-p}\norm{A}^{p}
        =n\bigl(1+\norm{n^{\mu-1}A}^{p}\bigr).
\]
The function $\mathcal L_p(Y)=\log(1+\norm{Y}^{p})$ is globally
$p$-Lipschitz on the underlying real Hilbert space.  Away from zero
its derivative has norm
$p\norm{Y}^{p-1}/(1+\norm{Y}^{p})\le p$, and the assertion at zero
follows by continuity.  Since $\beta_J$ and $\beta_J^0$ have the same
denominator, the decomposition \eqref{eq:src-18} therefore implies
\begin{equation}
  \log\beta_{J_j}-\log\beta_{J_j}^0
       \le Cp\,n_j^{-1+\mu}\norm{B_{J_j}}.
  \label{eq:src-27}
\end{equation}

\begin{lemma}[Square summation of residual blocks]
\label{res:block-square}
At a point $z\in Z$ where $R$ is a good path,
\begin{equation}
  \left(\sum_j\norm{B_{J_j}(z)}^2\right)^{1/2}\le L^{C_b}.
  \label{res:block-square-bound}
\end{equation}
Here the sum can be restricted to any subset of the possible blocks.
The fixed exponent $C_b$ is independent of $M$.
\end{lemma}

\begin{proof}
A residual packet in block $j$ has
\begin{equation}
  \frac{\rho_{j+1}}{B^2}<d_s\le\frac{\rho_j}{B^2},
  \qquad i_j\le k_s<i_{j+1}.
  \label{res:residual-range}
\end{equation}
Because $\rho_j/\rho_{j+1}\le B$, this block contains at most
$N_{\mathrm{len}}:=C(1+\log B)$ dyadic packet lengths.  Also $s_{T(s)},s_R\in J_j$,
so the top assignment and \eqref{res:residual-range} imply
\[
  \abs{\theta_s-s_R}
     \le C d_s+\rho_j
     \le C B^3 d_s\le C B^4d_s.
\]
Thus the reference-line estimates apply whenever $\ell_s\le l_R$.

The coefficient multiplying $q_{T(s)}\phi_s$ in $B_{J_j}$ is a
product of transition weights of modulus at most one, followed by
$n_{v_{k_s}(T(s))}^{-\mu/K}$.  The ending mass is at least
$h_{T(s)}$.  Consequently the modulus of the coefficient multiplying
$\phi_s$ is at most
$|q_{T(s)}|h_{T(s)}^{-\mu/K}\le C$.  This last bound follows from
the mask estimate \eqref{eq:src-8}: on the significant region
$h_T\ge1$, while off that region the mask has faster decay than
this fixed small negative power of $h_T$.

At $z$ let $e_\ell$ and $a_\ell$ be the sums of the absolute error and
main envelopes, respectively, at length $\ell$ in the fixed family
$\mathcal P_R$ of \eqref{hier:reference-packets}.  The preceding bounds
place every residual packet with $\ell_s\le l_R$ in this family, so its
absolute envelopes are included in these positive sums.  Goodness of $R$,
as specified in \Cref{hier:good-paths}, gives
\[
  \sum_{\ell\le l_R}e_\ell\le L^C,
  \qquad
  \left(\sum_{\ell\le l_R}a_\ell^2\right)^{1/2}\le L^C.
\]
The contribution of such lengths to $\norm{B_{J_j}}$ is bounded by
$C\sum_{\ell\in\mathcal I_j}(e_\ell+a_\ell)$, where
$\mathcal I_j$ is the set of residual lengths in that block.
The sets $\mathcal I_j$ are disjoint.  Positivity and
Cauchy--Schwarz within each block therefore give
\begin{align*}
 \left\{\sum_j
    \left(\sum_{\ell\in\mathcal I_j,\,\ell\le l_R}
                         e_\ell\right)^2\right\}^{1/2}
       &\le\sum_{\ell\le l_R}e_\ell,\\
 \sum_j\left(\sum_{\ell\in\mathcal I_j,\,\ell\le l_R}
                         a_\ell\right)^2
       &\le N_{\mathrm{len}}\sum_{\ell\le l_R}a_\ell^2.
\end{align*}
The factor $N_{\mathrm{len}}$ is polynomial in $L$.

Before the last possible block, $R$ remains eligible at the child
level.  Hence $d_R\le\rho_{j+1}/B^2<d_s$, and consequently
$\ell_s<l_R$.  Lengths larger than $l_R$ can occur only in the last
possible block.  At each of its at most $N_{\mathrm{len}}$ lengths, the total
absolute coefficient envelope is bounded by a constant, by
\Cref{conc:single-scale}, applied with reference length equal to
that scale.  To recall the reason this bound has no
factor counting the available labels, the condition
$|\theta_s-s_R|\le CB^4d_s$ makes the analyzing and synthesizing
box envelopes comparable to common $s_R$-coordinate envelopes, since
$W\gg B^4w$.  Summing centers first produces a common positive
convolution kernel of integral $O(1)$.  At each source point only
boundedly many selector labels are nonzero.  Summing these labels
then uses $\norm{g}_{\mathcal H}\le1$, and costs $O(1)$.
The uniformly bounded multipliers just considered preserve this
absolute estimate.  The last block therefore contributes at most
$C N_{\mathrm{len}}$ in norm.  Combining the estimates proves
\eqref{res:block-square-bound}.
\end{proof}

By \Cref{hier:marking}, at any marked point one has $h_j^*\ge c_m$
for a fixed $c_m>0$.  Monotonicity of the masses therefore implies
\begin{equation}
  \#\{j:\ j\text{ marked at }z,\ t\le n_j<2t\}
         \le C(1+t),\qquad t\ge M.
  \label{res:marked-shell-count}
\end{equation}
For completeness, all but the final marked index in this shell
account for disjoint mass drops of size at least $c_m$ while the mass
lies in an interval of length at most $t$; the final index costs one
even if its drop leaves the shell.  Dyadic summation gives
\begin{equation}
 \sum_{j\text{ marked at }z}n_j^{-2+2\mu}
       \le C\sum_{r\ge0}(2^rM)^{-1+2\mu}
       \le C M^{-1+2\mu}.
 \label{res:marked-weight-sum}
\end{equation}
Combining \eqref{eq:src-27}, \Cref{res:block-square}, and
Cauchy--Schwarz yields
\begin{equation}
 \sum_{j\text{ marked at }z}
       \bigl(\log\beta_{J_j}-\log\beta_{J_j}^0\bigr)
       \le CpL^{C_b}M^{-1/2+\mu}.
 \label{res:marked-cost}
\end{equation}

\subsection{Unary segments and their energy telescope}

Call an active index \emph{unary} at $z$ if its cell is unmarked and
$h_j^*(z)<1$.  At such an index the path child exists: at the last
possible block the drop would instead be $n_j>M$.  A significant top
has mass at least one.  Thus every significant top eligible in $J_j$
must remain eligible in the path child, since the total mass removed
from that child is $h_j^*<1$.

At a unary index write
\begin{equation}
  A_j=\lambda_j A_{j+1}+U_j+W_j.
  \label{res:unary-decomposition}
\end{equation}
Here $\lambda_j$ is the path-child coefficient in
\eqref{eq:src-18}; $U_j$ consists of the residual packets assigned
to tops which remain eligible in the path child.  These are exactly
the packets of the fixed subset $\mathcal Q_{R,\sigma}$ in
\eqref{hier:deterministic-residual-subset} which stop in block $j$.  The term $W_j$
collects all contributions of tops eligible in $J_j$ which do not
continue in the path child, including their residual packets and,
when present, their contributions through other children.

Those latter tops are all insignificant at $z$.  Their contributions
within $A_{J_j}$ are whole interval sums on individual tops, so
\Cref{hier:amplitudes} bounds the contribution of each by
$L^C(1+\norm{\zeta})|q_T|h_T^4$.  Since $h_T$ is bounded above by
an absolute constant and $|q_T|\le e^{-L\log L}$ there, this is at
most $Ce^{-L}h_T$ for sufficiently large $L$.  The masses of these
tops sum to $h_j^*$.  Therefore
\begin{equation}
  \norm{W_j}\le Ce^{-L}h_j^*.
  \label{res:discarded-top-error}
\end{equation}
This estimate uses a sum of top masses, not the number of such tops.

Partition each consecutive run of unary indices into consecutive
segments $[e,f)$.  Starting at $e$, end a completed segment as soon
as its accumulated drop is at least one.  Each individual drop is
less than one, so every completed segment has total drop in $[1,2)$.
At the end of the run, retain the final remainder as one segment if
it is nonempty.  Hence every segment satisfies
\begin{equation}
  \Delta_{e,f}:=\sum_{e\le j<f}h_j^*=n_e-n_f<2,
  \qquad n_f>M-2,
  \qquad n_j\asymp n_e\quad(e\le j\le f).
  \label{res:segment-masses}
\end{equation}
The endpoint remains an eligible path node, including when it is the
node where the mass stopping rule first applies.

The number of segment starts in a mass shell obeys
\begin{equation}
  \#\{[e,f):\ t\le n_e<2t\}\le C(1+t),\qquad t\ge M.
  \label{res:segment-shell-count}
\end{equation}
Completed segments account for disjoint drops of size at least one.
For each remainder other than the final remainder on the path,
charge it to the next interruption of a unary run.  Such an
interruption is either marked, with drop at least $c_m$, or is an
active unmarked index with drop at least one.  The mass at this
interruption differs from the start mass of the remainder by less
than two.  Thus the interruptions charged by starts in $[t,2t)$
lie in the enlarged shell $[t-2,2t)$ and their fixed positive drops
give $O(1+t)$ choices.  The one possible final remainder adds one.

For one segment define
\[
 A_i^*=\left(\prod_{i\le j<f}\lambda_j\right)A_f,
 \qquad V_i^*=V(n_i,A_i^*),\qquad e\le i\le f,
\]
with the empty product equal to one.  Since $n_i\ge n_{i+1}$,
$\alpha-p<0$, and $|\lambda_i|\le1$, we have
\begin{equation}
  V_i^*\le h_i^*+V_{i+1}^*,\qquad V_f^*=V_f.
  \label{res:starred-subsolution}
\end{equation}
Indeed, the baseline terms differ by $h_i^*$, while the energy
term on the left is at most
$n_{i+1}^{\alpha-p}\norm{A_{i+1}^*}^p$.
The denominator of $\beta_{J_i}$ contains the energy of the path
child and the nonnegative energies of all other children.  It is
therefore at least $h_i^*+V_{i+1}$.  It follows that
\[
 \sum_{e\le i<f}\log\beta_{J_i}
 \le \log V_e-\log V_f
       -\sum_{e\le i<f}\log(1+h_i^*/V_{i+1}).
\]
On the other hand, \eqref{res:starred-subsolution} gives
\[
 \log V_e^*-\log V_f
       \le\sum_{e\le i<f}\log(1+h_i^*/V_{i+1}^*).
\]
Subtracting these inequalities produces the useful estimate
\begin{equation}
 \begin{aligned}
 \sum_{e\le i<f}\log\beta_{J_i}
 &\le \log(V_e/V_e^*)
       +\sum_{e\le i<f}
           \left[\log(1+h_i^*/V_{i+1}^*)
                 -\log(1+h_i^*/V_{i+1})\right] \\
 &\le Cp\, n_e^{-1+\mu}
            \max_{e\le i\le f}\norm{A_i-A_i^*}.
 \end{aligned}
 \label{eq:src-28}
\end{equation}
To justify the last line, first use the Lipschitz estimate for
$\mathcal L_p$ to obtain
\[
  |\log V(n,A)-\log V(n,A')|
       \le p n^{\mu-1}\norm{A-A'}.
\]
For $h\ge0$, the derivative of $x\mapsto\log(1+he^{-x})$ has
absolute value $h/(e^x+h)$.  Along the interval between
$\log V_{i+1}$ and $\log V_{i+1}^*$, both endpoint amounts, and
therefore every intermediate amount, are at least $n_{i+1}$.
The corresponding correction term is bounded by
\[
 \frac{h_i^*}{n_{i+1}}
      p n_{i+1}^{\mu-1}\norm{A_{i+1}-A_{i+1}^*}.
\]
Now use \eqref{res:segment-masses} and
$\sum_{e\le i<f}h_i^*/n_{i+1}\le2/(M-2)$.  This proves
\eqref{eq:src-28}, including its endpoint case.

\subsection{Common residual weights and variation}

We give the coefficient calculation needed to apply the good-path
variation estimate.  It is important here that the same angular
grid underlies all shifts.

For integer levels $k$ along the path put
$v_k=v_k(R)$ and $a_k=n_{v_k}$, so $a_{i_j}=n_j$.
For a fixed segment $[e,f)$ and a block index $j\in[e,f)$, define
\begin{equation}
 \omega_j(k)=a_k^{-\mu/K}
     \prod_{r=i_j}^{k-1}
       \left(\frac{a_{r+1}}{m_{v_{r+1}}}\right)^{\zeta_{r\bmod K}},
 \qquad i_j\le k<i_f.
 \label{res:common-weight}
\end{equation}
Here $i_j$ denotes the angular level at the start of block $j$.

Every packet in $U_b$, for $j\le b<f$, is assigned to a top in
$v_{i_{b+1}}(R)$.  It follows by nesting that its top lies in $v_k(R)$
at every intermediate level $k\le i_{b+1}$, in particular at its
stopping level $k_s$.  Multiplying its coefficient in $U_b$ by the
path propagation factor $\prod_{j\le r<b}\lambda_r$ therefore
gives exactly $\omega_j(k_s)q_{T(s)}\phi_s$.
This identity holds for every contributing top, whether or not that
top is significant.  Thus the scalar coefficient depends only on
the stopping level, rather than on the packet or top inside that
level.

\begin{lemma}[Variation of the common coefficients]
\label{res:common-weight-variation}
Uniformly over the segments and their starting indices,
\begin{equation}
  \sup_{i_j\le k<i_f}|\omega_j(k)|
    +\sum_{k=i_j}^{i_f-2}|\omega_j(k+1)-\omega_j(k)|
       \le L^{C_w}.
  \label{res:common-weight-BV}
\end{equation}
The exponent $C_w$ is fixed before $M$ is chosen.
\end{lemma}

\begin{proof}
Every ending mass in this sequence is at least
$a_{i_f}=n_f>M-2$.  Since all real parts of $\zeta$ are nonnegative,
$|\omega_j(k)|\le a_k^{-\mu/K}\le1$ when $M\ge3$.
Using the real logarithms of positive masses, a logarithmic
increment of the coefficient is
\[
 \log\omega_j(k+1)-\log\omega_j(k)
 =\frac{\mu}{K}\log\frac{a_k}{a_{k+1}}
     -\zeta_{k\bmod K}\log\frac{m_{v_{k+1}}}{a_{k+1}}.
\]
The logarithm here is the specified linear expression defining the
coefficient, so no branch of a complex logarithm is being chosen.
The first sum of absolute increments telescopes.  For the second
use \Cref{hier:weight-variation} on the entire root-to-endpoint
path.  All its summands are nonnegative, whence
\begin{align*}
 \sum_{k=i_j}^{i_f-2}
    |\log\omega_j(k+1)-\log\omega_j(k)|
 &\le \left(\frac{\mu}{K}
          +K\max_\sigma|\zeta_\sigma|\right)
            \log\frac{n_{\mathrm{root}}}{n_f} \\
 &\le L^{C_w}.
\end{align*}
The last step uses $n_{\mathrm{root}}\le NL^C$ on $Z$,
$\log N\le2L$, $n_f> M-2\ge1$, and
$\max_\sigma|\zeta_\sigma|\le\mu+L^2$.
The full-root bound is needed: a denominator $m_v$ belonging to a
different shift may involve a block whose starting level precedes
the segment.  No estimate using only $\log(n_e/n_f)$ is asserted.

For two successive logarithmic values $z_0,z_1$ in this sequence,
both $|e^{z_0}|$ and $|e^{z_1}|$ are at most one.  Along the straight
segment between them, $|e^{(1-t)z_0+tz_1}|\le1$.  Integration of
the derivative of the exponential consequently yields
$|e^{z_1}-e^{z_0}|\le|z_1-z_0|$.
The logarithmic-increment bound proves
\eqref{res:common-weight-BV}.
\end{proof}

Let $\mathcal Q_{R,\sigma}$ denote the fixed residual subset in
\eqref{hier:deterministic-residual-subset}.  For each unary segment $[e,f)$ define
\[
 \Lambda_{e,f}(z)=
    \max_{i_e\le a\le b<i_f}
       \left\|\sum_{\substack{s\in\mathcal Q_{R,\sigma}\\
                              a\le k_s\le b}}
                       q_{T(s)}(z)\phi_s(z)\right\|.
\]
An empty sum has value zero.  Finite depth makes the maximum finite.
The interval of stopping levels $[a,b]$ corresponds to an interval
of dyadic packet lengths.  The latter follows directly from the
definition $k_s=\max\{k:d_s\le d(k)/B^2\}$, which is monotone
in $\ell_s$.  Since distinct unary segments have disjoint level
intervals $[i_e,i_f)$, one may choose a maximizing interval in each
segment and apply the single good-path $3$-variation bound to all
those disjoint intervals.  Thus
\begin{equation}
  \left(\sum_{[e,f)}\Lambda_{e,f}(z)^3\right)^{1/3}
       \le L^{C_v}.
  \label{res:segment-cube-sum}
\end{equation}
The segments may depend on $z$; the variation is a supremum over
all choices of disjoint intervals, so this dependence does not
alter the inequality.  The underlying packet subset itself is the
single deterministic set $\mathcal Q_{R,\sigma}$.

Abel summation for a finite scalar sequence $(\omega_k)$ and vector
sequence $(F_k)$ gives
\[
 \left\|\sum_{k=a}^b\omega_k F_k\right\|
 \le \left(|\omega_b|+
             \sum_{k=a}^{b-1}|\omega_k-\omega_{k+1}|\right)
                \max_{a\le t\le b}
                       \left\|\sum_{k=a}^tF_k\right\|.
\]
Applying this identity with \eqref{res:common-weight}, and then
iterating \eqref{res:unary-decomposition}, yields
\begin{equation}
 \max_{e\le i\le f}\norm{A_i-A_i^*}
    \le L^{C_w}\Lambda_{e,f}
           +Ce^{-L}\Delta_{e,f}.
 \label{res:segment-perturbation}
\end{equation}
Indeed the propagated $W_j$ terms have modulus at most their
unpropagated bounds because $|\lambda_j|\le1$.
There are no residual terms beyond $i_f$ in this calculation:
$U_j$ uses only $i_j\le k_s<i_{j+1}$ with $j<f$.
When $f$ is a mass-stopping node it is still an eligible node and
\eqref{res:segment-masses} applies.  A final block ending by
ineligibility cannot be unary and is excluded from this calculation.

Using \eqref{res:segment-shell-count} and dyadic mass shells, we get
\begin{equation}
 \begin{split}
 \sum_{[e,f)}n_e^{(-1+\mu)3/2}
   &\le C\sum_{r\ge0}(2^rM)^{-1/2+3\mu/2}
    \le C M^{-1/2+3\mu/2},\\
 \left(\sum_{[e,f)}n_e^{(-1+\mu)3/2}\right)^{2/3}
   &\le C M^{-1/3+\mu}.
 \end{split}
 \label{res:segment-weight-sum}
\end{equation}
H\"older with exponents $3/2$ and $3$, followed by
\eqref{eq:src-28}, \eqref{res:segment-cube-sum}, and
\eqref{res:segment-perturbation}, therefore bounds the sum of the
unary log factors by
\begin{equation}
 \sum_{j\text{ unary at }z}\log\beta_{J_j}
       \le CpL^{C_u}M^{-1/3+\mu}
          +Cp e^{-L}\left(\max_{j\text{ active}}n_j^\mu\right)
                              \sum_{j\text{ active}}u_j.
 \label{res:unary-cost}
\end{equation}
For the error term, use $n_e\ge n_j$ and $\mu-1<0$ within a
segment to obtain
$n_e^{\mu-1}h_j^*\le n_j^{\mu-1}h_j^*=n_j^\mu u_j$.
By \eqref{eq:src-12}, $\max n_j^\mu\le\exp(C\sqrt L)$.
Together with \eqref{res:relative-drop-sum}, this makes the second
term in \eqref{res:unary-cost} tend to zero.

Choose a fixed exponent $C_{\mathrm{res}}$ larger than all polynomial
exponents in \eqref{res:marked-cost} and \eqref{res:unary-cost},
and then choose
\begin{equation}
  A_M>4(C_{\mathrm{res}}+2),\qquad M=L^{A_M}.
  \label{res:choice-M}
\end{equation}
By \eqref{res:small-mu}, $-1/3+\mu\le-1/4$ and
$-1/2+\mu\le-5/12$.  Since $p=L^{1/2}$, this choice makes both
principal residual costs bounded by an absolute constant for
sufficiently large $L$.  None of the good-path thresholds or the
exponents used to select $A_M$ depended on $M$.

\subsection{The path product and its integral}

\begin{theorem}[Integrated path product]
\label{res:path-product}
With $M$ chosen by \eqref{res:choice-M}, the estimate
\eqref{eq:src-22} holds for every path top $R$, every shift, and
every vertex parameter with $\norm{\Im\zeta}_2\le L^2$.
Its constants are independent of the tile collection, angular
depth, and Hilbert-space dimension.
\end{theorem}

\begin{proof}
Fix a line of slope $s_R$, and on that line use the functions
$\Gamma_j$ constructed in \Cref{pred:marked-comparison}.
At a point of $Z$ where $R$ is good, the active indices fall into
the marked indices, the unary indices, and at most one active
unmarked index with $h_j^*\ge1$.
For the marked indices, combine \eqref{eq:src-24},
\eqref{res:marked-cost}, and
\eqref{res:relative-drop-sum}.  Their phase-domination losses have
sum at most $CLG^{-c}$, which is bounded.  The unary factors have
bounded log product by \eqref{res:unary-cost} and
\eqref{res:choice-M}.

At the possible remaining index, \Cref{hier:marking} gives
$M<n_j<M+1$.  From \eqref{eq:src-17} and the baseline $n_j$ in
the denominator of $\beta_{J_j}$,
\[
  \beta_{J_j}\le 1+L^{Cp}n_j^{\alpha-1}
       \le L^{Cp}(1+M)^C.
\]
Here $C$ is fixed, since $\alpha$ stays in a fixed compact interval
below three.  We have therefore proved the pointwise domination
\begin{equation}
 \prod_{j\text{ active at }z}\beta_{J_j}(z)
       \le L^{Cp}(1+M)^C\prod_{j=0}^{J_R}\Gamma_j(z),
 \qquad z\in Z\cap\{R\text{ good}\},
 \label{res:pointwise-product}
\end{equation}
where $J_R$ is the last possible block index for $R$.
It is legitimate to include every possible index on the right:
if $n_j(z)\le M$, its mesh cell cannot satisfy
$\inf n_j>M$, so it is unmarked and $\Gamma_j(z)=1$.
Unmarked active indices likewise have $\Gamma_j=1$.

Recall that $\Gamma_j$ is constant on the mesh of length $H_{j-1}$
and has average at most one on every cell of length $H_j$.
The meshes are nested.  On any cell $I$ of the largest mesh
$H_{J_R}$, successive averaging from the smallest mesh upward
therefore gives
\begin{equation}
  \int_I\prod_{j=0}^{J_R}\Gamma_j(x)\,\dd x\le |I|.
  \label{res:mesh-average}
\end{equation}
More explicitly, all later factors $\Gamma_k$, $k>j$, are constant
on each $H_j$ cell.  Integrating $\Gamma_j$ on those cells can thus
only decrease the integral of the nonnegative product.  Iterate
this observation for $j=0,1,\ldots,J_R$.

The weight $h_R$ is uniformly comparable on $I$.  This remains true
for a terminal block and for the short first block: eligibility
of $R$ at $i_j$ gives $\rho_j\ge B^2d_R$, and every block has
$\rho_j/\rho_{j+1}\le B$.  Thus
\begin{equation}
  H_j=\frac{w}{\rho_{j+1}\sqrt G}
        \le\frac{l_R}{B\sqrt G}\le l_R.
  \label{res:mesh-versus-top}
\end{equation}
Along the line of slope $s_R$, the transverse normalized coordinate
of $R$ is fixed and the horizontal normalized coordinate changes
by at most $H_j/l_R$.  The definition \eqref{eq:src-5} consequently
gives $\sup_I h_R\le C\inf_Ih_R$, with an absolute constant
depending only on the fixed decay exponent.  By
\eqref{res:mesh-average},
\[
 \int_I h_R\prod_{j=0}^{J_R}\Gamma_j\,\dd x
     \le (\sup_Ih_R)|I|\le C\int_Ih_R\,\dd x.
\]
Sum over the largest-mesh cells and use
\eqref{res:pointwise-product}.  Whenever the original integrand in
\eqref{eq:src-22} has a measurable slice on this line, we obtain
\[
 \int_{Z\cap\{R\text{ good}\}\text{ on the line}}
       h_R\prod_{j\text{ active}}\beta_{J_j}\,\dd x
       \le L^{Cp}(1+M)^C
                    \int_{\text{line}}h_R\,\dd x.
\]

The snapshots used to construct the $\Gamma_j$ need only be chosen
on this one line.  There are countably many mesh cells, so these
choices give measurable step functions on that line.  No jointly
measurable choice of snapshots across different lines is required.
The original integrand in \eqref{eq:src-22}, expressed entirely through
$h_R$, the good sets, and the finite hierarchy factors $\beta_{J_j}$, is
nonnegative and Lebesgue measurable.  In the coordinates $(x,y-s_Rx)$,
it has a measurable slice for almost every transverse parameter.  The
preceding line estimate therefore holds for almost every such parameter.
Integrating it and applying Tonelli, using the Jacobian-one coordinate
change, proves
\eqref{eq:src-22}.
\end{proof}

All sums and products in this section were finite.  The shell counts,
Abel estimates, and mesh averaging bounds contain no factor equal to
the number of scales, tiles, or hierarchy levels.  They therefore
apply with identical constants to every finite truncation required
in the packet reduction.
